\documentclass[book.tex]{subfiles}
\begin{document}
\chapter{Molecular Representations and Featurizations}
\label{chap:molecular_ml}
\noindent\rule{11cm}{0.4pt} \\
\textbf{Prerequisites:}
\autoref{chap:graphconvs}, \autoref{chap:molecular_hamiltonian} \\
\textbf{Difficulty Level:} ** \\
\noindent\rule{11cm}{0.4pt}
\vspace{5mm}

Molecules, groups of two or more atoms held together by chemical bonds, form the smallest repeating units of materials. Identifying the optimal molecule for an application, such as designing a medical drug or a battery electrolyte, depends on emergent properties that emerge from interactions at the molecular length scale. Models that predict emergent collective properties are consequently of enormous value in practical applications. In this chapter, we will learn about the basics of molecular machine learning and molecular featurization. Molecular featurizations use heuristic methods to compute useful properties of a molecular structure for use in downstream learning pipelines. 

\section{Molecular Featurization}

In order to use machine learning to predict molecular properties, we need to encode properties of the molecule, or "features" that a learning algorithm can discern patterns from. Several molecular featurization techniques have been proposed in order to calculate molecular descriptors such as functional groups, residual charges, dipole moments, inter-atom distance metrics and bonding features. Especially in more complex systems, generating refined molecular features becomes important to capture the subtle interplay of interactions such as hydrogen bonding and van del Waals forces in predicting collective behavior and emergent properties. Generated features are typically used as input to train appropriate machine learning models to predict collective properties that are relevant in the material selection process. 
%methods.\cite{gomez2018automatic} 
\subsection{Smiles Representation of Molecules}

Molecular machine learning relies on effective encoding of molecules into fixed-length strings or vectors. The simplified molecular-input line-entry system (SMILES) refers to a line notation (a typographical method using printable characters) for encoding the composition, bonding and structure of molecules. 

Examples of molecules and their associated SMILES stirngs are given in Table~\ref{table:smiles}. SMILES strings can be used to uniquely represent molecules, and enable an effective starting point for molecular featurization methods. Such methods process the SMILES to generate fixed-length feature vectors and/or matrices. Note that the SMILES representation includes only heavy (non hydrogen) atoms with the assumption that unsatisfied bonds are occupied by hydrogen atoms
%\href{https://www.daylight.com/dayhtml/doc/theory/theory.smiles.html}{SMILES specification rules} are important for consistency and uniqueness of representation.  A few examples are shown below:
\begin{table}[!ht]
\centering
\begin{tabular}{|c |c| c|}
\hline
Molecule & Molecular Formula & SMILES string/ID \\
\hline
Oxygen & O$_2$ & O=O \\
Propane & CH$_3$CH$_2$CH$_3$& CCC \\
Ethanol & CH$_3$CH$_2$OH & CCO \\
Acetic acid & CH$_3$COOH& CC(=O)O\\ 
Cyclohexane & C$_6$H$_{12}$&C1CCCCC1\\
Hydrogen Cyanide & HCN & C\#N\\
\hline
\end{tabular}
\label{table:smiles}
\end{table}

SMILES strings offer a way to uniquely represent molecules as strings, but most molecular machine learning methods require additional information such as bonding characteristics, composition and connectivity to learn sophisticated features of molecules from limited amounts of data. Prior work has demonstrated that SMILES strings can be used to predict emergent properties based on the ability to effectively encode molecules from SMILES representations using more sophisticated representations.

%A few useful molecular featurization methods (also implemented in MoleculeNet) that can operate on molecular representations such as the SMILES are discussed below.
\subsection{Extended Connectivity Fingerprints (ECFP)}
Extended-Connectivity fingerprints are extensively used for molecular characterizations in chemical informatics\cite{rogers2010extended}. These fingerprints transform a SMILES representation of a molecule into a sparse bitvector that can be used in a downstream machine learning pipeline. To generate the ECFP vector, a molecule is first decomposed into submodules originated from heavy atoms (non-hydrogen), each assigned with a unique identifier. The submodules (segments) and identifiers are extended through bonds to generate larger substructures and associated identifiers. These substructures are hashed into a fixed length binary fingerprint. This approach to encoding molecules enables it to be applied to tasks such as similarity searching and prediction of collective properties since it contains information about topological characteristics of the molecule.

To implement the circular fingerprint algorithm, we will assume we have access to some hash function with the following type signature
\begin{lstlisting}[language=python]
hash: $\forall$ n, $\mathbb{Z}_2^N$ -> $\mathbb{Z}_2^k$
\end{lstlisting}
We then define the fingerprint computation itself below, where \lstinline{R} is the neighbor radius considered, \lstinline{S} is the size of the fingerprint, and \lstinline{g} is a mapping of each atom to a bit vector of features.
%\textcolor{red}{TODO: Make Proper Physika}
\begin{lstlisting}[language=python]
class CircularFingerprint(R: $\mathbb{N}$, S: $\mathbb{N}$, g: Node -> $\mathbb{Z}_2^k$):

  def $\lambda$(molecule: Graph):
    f = zeros(S)
    # Loop over atoms
    for a in molecule: 
      r[a] = g(a)
    for L in [1,$\dotsc$, R]:
      nbrs = neighbors(a)
      vec = r[a] || nbrs
      r[a] = hash(vec)
      i = mod(r[a], S)
      f[i] = 1
    return f
\end{lstlisting}

\section{The Design Workflow}

A typical design process for identifying molecular candidates for any given application involves the following steps:

\begin{itemize}
    \item Identifying quantitative design criteria and associated molecular properties of relevance, and using these to generate features summarizing the molecule.
    \item Training an appropriate machine learning model (performing hyperparameter and architecture optimization, etc.), and then validating and testing the model's ability to predict the identified molecular properties of interest.
    \item Identifying candidates that meet the quantitative design criteria using predictions from machine learning models 
    \item Performing experimental testing to measure the true properties of promising candidates, and creating a feedback loop with new training data to refine the featurization and prediction models 
    %\item Identifying a set of improved candidates for the application
\end{itemize}

This design loop is often called the active learning process whereby machine learning models are used to guide experimental design. We discuss it here in the context of molecules, but more generally a similar process applies for materials discovery or engineering design as well.

\begin{figure}[ht]
   % \includegraphics[width=\linewidth]{figures/Differentiable Physics/Molecular ML/Molecular Featurizations/typicalWorkflow.png}
   \includegraphics[width=\linewidth]{figures/Differentiable Physics/Molecular ML/Molecular Featurizations/ML_typical_Workflow.png}
   \caption{Typical workflow for making machine learning based property predictions for molecular structures. A molecular structure is transformed into a numerical representation called features or descriptors. The features are then used as an input for a machine learning model that is trained to output a property for the structure. There is also a possibility of combining the descriptor and learning model together into one inseparable step. %\textcolor{red}{Redo this figure and include a feedback loop-- source  http://dx.doi.org/10.17632/vzrs8n8pk6.1}
   }
\label{fig:typicalWorkflow}
\end{figure}

In subsequent chapters, we will expand on different choices for each of these modeling steps.

\section{Exercises}
\begin{enumerate}
    \item Generate the SMILES string for acetaminophen (HOC$_6$H$_4$NHCOCH$_3$)
    \item Compute the ECFP fingerprint for this molecule.
\end{enumerate}


\end{document}
